\clearpage
\begin{trapbox}
\textbf{\textcolor{CTrap}{易错一｜把“加 $a$ 减 $a$”误记成“周期 $a$”}}
\begin{center}\begin{tikzpicture}[font=\small,>=stealth,line join=round]
\draw[gray!40,rounded corners=3pt] (0,0) rectangle (14.2,8.3);\node at (7.1,7.700000000000001) {算的是两个输入之差，不是中点到单侧的距离};
\draw[->,gray] (.8,3.2)--(13.1,3.2) node[right]{$x$};\foreach \v in {1,3,5,7,9}{\pgfmathsetmacro{\X}{1.8+\v}\draw[gray] (\X,3.2)--(\X,3.1);\node[below=7pt] at (\X,3.1){$\v$};}\foreach \v in {1,5,9}{\fill[blue] ({1.8+\v},3.2) circle(2.4pt);}\draw[<->,orange!85!black,thick] (2.8,4.7)--(6.8,4.7) node[midway,above]{$5-1=4$};\draw[<->,violet,thick] (6.8,6.1)--(10.8,6.1) node[midway,above]{$9-5=4$};\node at (7.1,1.8){$x=3$ 时，输入为 $x-2=1$ 与 $x+2=5$；中点到两边各2。};
\node[text width=12.6cm,align=center] at (7.1,.65) {相邻标点相隔4；仅凭这几个点，不能证明整个函数的周期。};
\end{tikzpicture}\end{center}



例如 $f(x+2)=f(x-2)$，比较的其实是 $x+2$ 和 $x-2$，相隔 $4$，并非 $2$。因此直接得到的是\textbf{一个正周期 $4$}。

\medskip

\end{trapbox}
\clearpage
\begin{trapbox}
\textbf{\textcolor{CTrap}{易错二｜把“一个周期”写成“最小正周期”}}
\begin{center}\begin{tikzpicture}[font=\small,>=stealth,line join=round]
\draw[gray!40,rounded corners=3pt] (0,0) rectangle (14.2,8.3);\node at (7.1,7.700000000000001) {例6：条件保证周期4，曲线却每隔2已经重复};
\begin{scope}[shift={(1.6882352941176473,3.7409090909090907)},x=1.1764705882352942cm,y=1.3939393939393938cm]\draw[->,gray] (-0.5,0)--(4.6,0) node[right]{$x$};\draw[->,gray] (0,-1.5)--(0,1.8) node[above]{$y$};\draw[blue,thick,solid,domain=-0.2:4.3,samples=180] plot(\x,{cos(deg(pi*\x))});\fill[orange!85!black] (0,1) circle(2.2pt);\fill[orange!85!black] (2,1) circle(2.2pt);\fill[orange!85!black] (4,1) circle(2.2pt);\end{scope}\node[anchor=west,align=left,text width=5.0cm] at (8.4,4.0) {$f=\cos(\pi x)$\\[9pt]$f(x+2)=f(x-2)$\\[9pt]推出4是一个周期\\[9pt]最小正周期为2\\[9pt]图中4含两个完整重复};
\node[text width=12.6cm,align=center] at (7.1,.65) {只有进一步使用这条具体函数的性质，才能确定最小正周期。};
\end{tikzpicture}\end{center}



$f(x+2)=f(x-2)$ 只能保证 $4$ 是一个周期。反例：$f(x)=\cos(\pi x)$ 的最小正周期是 $2$，但也满足该条件。甚至常数函数满足条件，却没有最小正周期。

\textbf{高考用语：} 如果没有进一步信息，写“$T=4$ 是一个正周期”；\textbf{不要直接写“最小正周期为 $4$”。}

\medskip

\end{trapbox}
\clearpage
\begin{trapbox}
\textbf{\textcolor{CTrap}{易错三｜忘记 $a$ 可以是负数}}
\begin{center}\begin{tikzpicture}[font=\small,>=stealth,line join=round]
\draw[gray!40,rounded corners=3pt] (0,0) rectangle (14.2,8.3);\node at (7.1,7.700000000000001) {例5：负数表示走向左，正周期仍取距离6};
\begin{scope}[shift={(4.1,3.7409090909090907)},x=0.42857142857142855cm,y=1.3939393939393938cm]\draw[->,gray] (-7,0)--(7,0) node[right]{$x$};\draw[->,gray] (0,-1.5)--(0,1.8) node[above]{$y$};\draw[blue,thick,solid,domain=-6.5:6.5,samples=180] plot(\x,{cos(deg(pi*\x/3))});\fill[orange!85!black] (-3,-1) circle(2.2pt);\fill[orange!85!black] (3,-1) circle(2.2pt);\fill[orange!85!black] (0,1) circle(2.2pt);\fill[orange!85!black] (6,1) circle(2.2pt);\end{scope}\node[anchor=west,align=left,text width=5.0cm] at (8.4,4.0) {示例 $f=\cos\frac{\pi x}3$\\[9pt]$f(x-3)=f(x+3)$\\[9pt]取 $a=-3$，$2a=-6$\\[9pt]移动方向可以向左\\[9pt]正周期 $T=6$};
\node[text width=12.6cm,align=center] at (7.1,.65) {正周期写非零移动量的绝对值；不是把带符号的位移直接当成正数长度。};
\end{tikzpicture}\end{center}



给出 $f(x-3)=f(x+3)$，从等式可推出带符号的移动量 $-6$；如果要求\textbf{正周期}，必须写 $T=6$，即一般写 $T=2|a|$。

\medskip

\end{trapbox}
\clearpage
\begin{trapbox}
\textbf{\textcolor{CTrap}{易错四｜一个或几个相等的函数值，不能推出周期}}
\begin{center}\begin{tikzpicture}[font=\small,>=stealth,line join=round]
\draw[gray!40,rounded corners=3pt] (0,0) rectangle (14.2,8.3);\node at (7.1,7.700000000000001) {两点同高，只说明这两点：平方函数不周期};
\begin{scope}[shift={(3.1000000000000005,1.946774193548387)},x=0.7142857142857144cm,y=0.14838709677419354cm]\draw[->,gray] (-2.8,0)--(5.6,0) node[right]{$x$};\draw[->,gray] (0,-2)--(0,29) node[above]{$y$};\draw[blue,thick,solid,domain=-2.2:5,samples=180] plot(\x,{(\x)^2});\fill[orange!85!black] (-2,4) circle(2.2pt);\fill[orange!85!black] (2,4) circle(2.2pt);\fill[orange!85!black] (1,1) circle(2.2pt);\fill[orange!85!black] (5,25) circle(2.2pt);\end{scope}\node[anchor=west,align=left,text width=5.0cm] at (8.4,4.0) {$f=x^2$\\[9pt]$f(-2)=f(2)=4$\\[9pt]虽然两输入相隔4\\[9pt]但 $f(1)=1,\ f(5)=25$\\[9pt]不能处处平移4};
\node[text width=12.6cm,align=center] at (7.1,.65) {否定周期只需找到一对平移后不等的值；证明周期必须对所有输入成立。};
\end{tikzpicture}\end{center}



比如 $f(x)=x^2$ 时 $f(-2)=f(2)=4$，但 $x^2$ 并不是周期函数。\textbf{两个位置高度一样，只能说明这两点一样高；必须对所有 $x$ 都能按同一个固定间隔重复，才叫周期。}

\medskip

\end{trapbox}
\clearpage
\begin{trapbox}
\textbf{\textcolor{CTrap}{易错五｜“足够长的区间”不等于完整的周期定义域}}
\begin{center}\begin{tikzpicture}[font=\small,>=stealth,line join=round]
\draw[gray!40,rounded corners=3pt] (0,0) rectangle (14.2,8.3);\node at (7.1,7.700000000000001) {只截取两次完整波形，定义域仍不是整个实数轴};
\begin{scope}[shift={(1.6,3.7409090909090907)},x=0.5cm,y=1.3939393939393938cm]\draw[->,gray] (-1,0)--(11,0) node[right]{$x$};\draw[->,gray] (0,-1.5)--(0,1.8) node[above]{$y$};\draw[blue,thick,solid,domain=0:8,samples=180] plot(\x,{sin(deg(pi*\x/2))});\fill[orange!85!black] (0,0) circle(2.2pt);\fill[orange!85!black] (8,0) circle(2.2pt);\end{scope}\node[anchor=west,align=left,text width=5.0cm] at (8.4,4.0) {$D=[0,8]$\\[9pt]图取 $f=\sin\frac{\pi x}2$\\[9pt]可见两个相似波形\\[9pt]但 $x=7$ 加4后\\[4pt]就没有定义};
\node[text width=12.6cm,align=center] at (7.1,.65) {有界定义域无法沿非零周期来回平移；截取一段后须重新判断。};
\end{tikzpicture}\end{center}



主结论假定函数定义在整个 $\mathbb R$。如果只在一个有限区间里观察到相同高度，不能直接宣布函数在整个数轴上周期性重复；题目若只给某一段的条件，必须重新检查区间外是否也有定义、是否满足等式。

\medskip

\end{trapbox}
\clearpage
\begin{trapbox}
\textbf{\textcolor{CTrap}{易错六｜把对称轴条件和周期条件混为一谈}}
\begin{center}\begin{tikzpicture}[font=\small,>=stealth,line join=round]
\draw[gray!40,rounded corners=3pt] (0,0) rectangle (14.2,8.5);\node at (7.1,7.9) {固定中点的等高是轴对称；比较中点变化才对应此处周期};
\node at (2.4,6.9) {轴对称：$f=(x-2)^2$};\begin{scope}[shift={(1.5499999999999998,3.628571428571429)},x=0.425cm,y=0.5142857142857143cm]\draw[->,gray] (-2,0)--(6,0) node[right]{$x$};\draw[->,gray] (0,-2)--(0,5) node[above]{$y$};\draw[blue,thick,solid,domain=0:4,samples=160] plot(\x,{(\x-2)^2});\fill[orange!85!black] (1,1) circle(2pt);\fill[orange!85!black] (3,1) circle(2pt);\end{scope}\node[align=center,text width=4cm] at (2.4,1.7) {固定中点2：\\$f(1)=f(3)$};\node at (7.050000000000001,6.9) {周期示例：$g=\sin\frac{\pi x}2$};\begin{scope}[shift={(6.2,3.628571428571429)},x=0.425cm,y=0.5142857142857143cm]\draw[->,gray] (-2,0)--(6,0) node[right]{$x$};\draw[->,gray] (0,-2)--(0,5) node[above]{$y$};\draw[violet,thick,solid,domain=-1.8:5.8,samples=160] plot(\x,{sin(deg(pi*\x/2))});\fill[orange!85!black] (-1.5,-0.7071067811865476) circle(2pt);\fill[orange!85!black] (2.5,-0.7071067811865476) circle(2pt);\end{scope}\node[align=center,text width=4cm] at (7.050000000000001,1.7) {中点0.5：\\$g(-1.5)=g(2.5)$};\node at (11.7,6.9) {同一周期函数 $g$};\begin{scope}[shift={(10.85,3.628571428571429)},x=0.425cm,y=0.5142857142857143cm]\draw[->,gray] (-2,0)--(6,0) node[right]{$x$};\draw[->,gray] (0,-2)--(0,5) node[above]{$y$};\draw[violet,thick,solid,domain=-1.8:5.8,samples=160] plot(\x,{sin(deg(pi*\x/2))});\fill[orange!85!black] (0.5,0.7071067811865476) circle(2pt);\fill[orange!85!black] (4.5,0.7071067811865476) circle(2pt);\end{scope}\node[align=center,text width=4cm] at (11.7,1.7) {中点2.5：\\$g(0.5)=g(4.5)$};
\node[text width=12.6cm,align=center] at (7.1,.65) {三图同尺度；右侧两例的比较中点并非对称轴，配对点之间始终相距4。};
\end{tikzpicture}\end{center}



\begin{itemize}[leftmargin=*,itemsep=.3em]
\item $f(x+a)=f(x-a)$：两端围绕\textbf{不断变化的} $x$，当条件对所有 $x$ 成立时，推出周期。
\item $f(c+t)=f(c-t)$：两端围绕\textbf{固定的} $c$，表示图像关于 $x=c$ 对称。
\end{itemize}
特别注意：\textbf{这里说的是比较点的“中间位置”是否固定}，不能只看式子里有没有“加减同一个数”。

\medskip

\end{trapbox}
\clearpage
\begin{trapbox}
\textbf{\textcolor{CTrap}{易错七｜大数求值算错落点}}
\begin{center}\begin{tikzpicture}[font=\small,>=stealth,line join=round]
\draw[gray!40,rounded corners=3pt] (0,0) rectangle (14.2,8.3);\node at (7.1,7.700000000000001) {例3：先把大数移回基本区间，再代入平方};
\begin{scope}[shift={(2.9000000000000004,2.033333333333333)},x=0.4cm,y=0.3833333333333333cm]\draw[->,gray] (-4.5,0)--(10.5,0) node[right]{$x$};\draw[->,gray] (0,-1)--(0,11) node[above]{$y$};\draw[blue,thick,solid,domain=-4:-3,samples=180] plot(\x,{(\x+6)^2});\draw[blue,thick,solid,domain=-3:3,samples=180] plot(\x,{(\x)^2});\draw[blue,thick,solid,domain=3:9,samples=180] plot(\x,{(\x-6)^2});\draw[blue,thick,solid,domain=9:10,samples=180] plot(\x,{(\x-12)^2});\fill[orange!85!black] (2,4) circle(2.2pt);\fill[orange!85!black] (8,4) circle(2.2pt);\fill[orange!85!black] (-3,9) circle(2.2pt);\fill[orange!85!black] (3,9) circle(2.2pt);\fill[orange!85!black] (9,9) circle(2.2pt);\end{scope}\node[anchor=west,align=left,text width=5.0cm] at (8.4,4.0) {基本段 $[-3,3)$ 上 $f=x^2$\\[9pt]按周期6延拓\\[9pt]$2024=337\times6+2$\\[9pt]$f(2024)=f(2)=4$\\[9pt]图中 $f(2)=f(8)=4$};
\node[text width=12.6cm,align=center] at (7.1,.65) {图中每段来自同一基本段的平移；接头处实际取值为9，两侧也趋近9。};
\end{tikzpicture}\end{center}



若正周期为 $6$，要求 $f(2024)$，可以用 $2024=337\times6+2$ 得到 $f(2)$；但若已知解析式只适用于另一个区间，需继续加减 $6$，直到落在\textbf{给出具体表达式的范围内}，不能机械地总取余数。
\end{trapbox}
